\Needspace{6\baselineskip}
\begin{center}\hypertarget{native-250-article178}{}\pdfbookmark[2]{178.}{article-178}178.\end{center}
The function just obtained certainly cannot express the probabilities of errors with complete rigor. Since possible errors are always confined within definite bounds, the probability of larger errors ought always to become $=0$, whereas our formula always gives a finite value. This defect, which every analytic function must have by its nature, is of no consequence for any practical application. The value of our function decreases so rapidly once $h\Delta$ has attained a substantial value that it can safely be regarded as equivalent to $0$. Moreover, the nature of the problem will never permit the limits of the errors themselves to be assigned with absolute rigor.

\SourceLeaf{251}The constant $h$ may also be regarded as a measure of the precision of the observations. Suppose the probability of an error $\Delta$ is expressed in one set of observations by $\frac{h}{\BareRoot{\pi}}e^{-hh\Delta\Delta}$ and in another set of more or less precise observations by $\frac{h'}{\BareRoot{\pi}}e^{-h'h'\Delta\Delta}$. The expectation that the error in an observation of the first set lies between the limits $-\delta$ and $+\delta$ is expressed by the integral $\int\frac{h}{\BareRoot{\pi}}e^{-hh\Delta\Delta}\,\mathrm{d}\Delta$, taken from $\Delta=-\delta$ to $\Delta=+\delta$. Similarly, the expectation that the error in an observation of the second set remains within the limits $-\delta'$ and $+\delta'$ is expressed by the integral $\int\frac{h'}{\BareRoot{\pi}}e^{-h'h'\Delta\Delta}\,\mathrm{d}\Delta$, taken from $\Delta=-\delta'$ to $\Delta=+\delta'$. The two integrals are clearly equal whenever $h\delta=h'\delta'$. If, for example, $h'=2h$, an error of twice the size in the first set can occur just as readily as an error of half that size in the second. In ordinary usage the observations of the second set are then said to have twice the precision.
